\documentclass{article}
\usepackage{graphicx} % Required for inserting images
\usepackage{datetime} % Add this line to include the datetime package
\usepackage{amsmath}
\usepackage{amssymb}
\usepackage{amsthm}
\usepackage{tikz}
%\usepackage{hyperref}
\usepackage{xcolor} % Required for defining custom colors

\definecolor{darkblue}{rgb}{0.0, 0.0, 0.5} % Define a darker blue color

\usepackage[colorlinks=true, linkcolor=darkblue, urlcolor=darkblue, citecolor=darkblue]{hyperref} % Use the darker blue color

% Define theorem style
\theoremstyle{plain} % Bold title, italic body
\newtheorem{theorem}{Theorem}
\newtheorem{lemma}[theorem]{Lemma}
\newtheorem{corollary}[theorem]{Corollary}

\theoremstyle{definition} % Bold title, normal body
\newtheorem{definition}[theorem]{Definition}
\newtheorem{example}[theorem]{Example}


\title{Homework 4}
\author{Aaron Honculada}
\date{\monthname[\month] \the\year}

\begin{document}

\maketitle

\section*{Problem 1}
\noindent\textit{Show that $a^{1729} \equiv a \pmod{1729}$ for all integers $a$. (Note: $1729 = 7 \cdot 13 \cdot 19$.)}

\begin{proof}
    By Fermat's Little Theorem, we know that for primes, 7, 13 and 19,
    \begin{align*}
        a^6 &\equiv 1 \pmod{7} \\
        a^{12} &\equiv 1 \pmod{13} \\
        a^{18} &\equiv 1 \pmod{19}. \\
    \end{align*}
    We can also factor $1728 = 6 \cdot 288 = 12 \cdot 144 = 18 \cdot 96$.
    Which means we have:
    \begin{align*}
        a^{1728} = a^{(6)^{288}} = 1^{288} &\equiv 1 \pmod{7} \\
        a^{1728} = a^{(12)^{144}} = 1^{144} &\equiv 1 \pmod{13} \\
        a^{1728} = a^{(18)^{96}} = 1^{96} &\equiv 1 \pmod{19}. \\
    \end{align*}
    
    Given that 7, 13, and 19 are coprime, we can apply the Chinese Remainder Theorem
    which states that if:

    \begin{align*}
        a^{1729} &\equiv 1 \pmod{7} \\
        a^{1729} &\equiv 1 \pmod{13} \\
        a^{1729} &\equiv 1 \pmod{19} \\
    \end{align*}
    
    then:
    \begin{centering}
        $a^{1729} \equiv a \pmod{1729}$
    \end{centering}


\end{proof}

\section*{Problem 2}
\noindent\textit{For each example below, determine where the subgroup $H$ is a normal
subgroup of $G$.}

\noindent\text{(a) $G=\{\text{ functions }f: \mathbb{R} \to \mathbb{R}\}$ under addition, and
$H=\{f \mid f(1)=0\}$}

Given that the addition of two real valued functions is commutitive, i.e. for real valued functions
$f, f' \in G$ we have $f + f' = f' + f$, we know that $G$ is abelian, and every subgroup of $G$ is a normal subgroup.
Given two real valued functions $f,g \in G$, $f$ and $g$ belong to the same coset if
$(f-g)(1) = 0$ and since every $f(1) = r$, then each coset corresponds to an $r \in \mathbb{R}$.

Therefore, $G/H \cong \mathbb{R}$.

\noindent\text{(b) $G=GL(2,\mathbb{R})$, and $H= \left\{ \left(\begin{matrix}a & 0 \\ 0 & b \end{matrix}\right) \mid a, b \neq 0 \right\}$.}

Let $h= \left( \begin{matrix}
    a & 0 \\ 0 & b
\end{matrix} \right)$,
$g = \left( \begin{matrix}
    p & q \\ r & s
\end{matrix} \right)$, and 
$g^{-1} = \frac{1}{sp-qr}\left( \begin{matrix}
    s & -q \\ -r & p
\end{matrix} \right)$.

We evaluate $g^{-1}hg = $
$\frac{1}{sp-qr}\left( \begin{matrix}
    s & -q \\ -r & p
\end{matrix} \right)
\left( \begin{matrix}
    a & 0 \\ 0 & b
\end{matrix} \right)
\left( \begin{matrix}
    p & q \\ r & s
\end{matrix} \right)$.

$hg = \left( \begin{matrix}
    ap & aq \\ br & bs
\end{matrix} \right)$,
$g^{-1}hg = \frac{1}{ps-qr}\left( \begin{matrix}
    sap-qbr & saq - qbs \\ pbr - rap & pbs - raq
\end{matrix} \right)$.

The resulting matrix is not a diagonal matrix and thus $g^{-1}hg \not\in H$ and $H$ is not a normal subgroup of $G$.

\noindent\text{(c) $G=S_4$ and $H=\{e, (1\text{ }2)(3\text{ }4), (1\text{ }3)(2\text{ }4), (1\text{ }4)(2\text{ }3) \}$.}
We observe that $H$ is a finite group of 4 elements where each element is its own inverse. Thus $H \cong \text{Klein-four group}$.
We show that the Klein-four group is a normal subgroup of $S_4$. We define an surjective map $f: S_4/H \to S_3$ where the elements
of $H$ are mapped to the permutations of $\{1,2,3\}$, where the transpositions
containing 4 will fix the other element. For example, the permutation $(1\text{ }2)(3\text{ }4)$
fixes 3 and switches the order of 1 and 2 resulting in $(1\text{ }2)$.

$S_4/H \cong S_3$.

\section*{Problem 3}
\noindent\textit{The quaternion group $Q_8$, consists of eight elements \{1, -1, i, -i, j, -j, k, -k \}
such that $i^2=j^2=k^2=-1$.}

\noindent\text{(a) Show that $Q_8$ is not isomorphic to $D_4$.}

We see that among the elements of $Q_8$, the identity has order 1, the elements $-1,-i,-j,-k$ have order 2,
and the remaining elements $i, -i, j, -j, k, -k$ have order 4. Considering $D_4$ there are five elements of order 2, namely the 
four reflections and the rotation of 180 degrees. There is no way to construct an injective function sending all the elements
of order 2 in $D_4$ to the elements of order 2 in $Q_8$.
Thus, $Q_8 \not\cong D_4$.

\noindent\text{(b) Find all subgroups of $Q_8$.}

We list all subgroups of $Q_8$. 
\begin{itemize}
    \item $\langle 1 \rangle$
    \item $\{1, -1\}$
    \item $\{1, -1, i, -i\}$
    \item $\{1, -1, j, -j\}$
    \item $\{1, -1, k, -k\}$
    \item $Q_8$
\end{itemize}
The list is exhaustive due to Langrange's Theorem which implies that the order of
each subgroup must divide the order of $Q_8$. The multiples of 8 are 1, 2, 4, and 8.
Only the trivial subgroup is order 1 and $Q_8$ is the only subgroup of order 8, thus the any proper, non-trivial
subgroup must either be of order 2 or order 4. When we attempt to include an element to a subgroup,
we must also include its inverse, thus the order will always increase by 2. The subgroup $\langle -1 \rangle$ is the
only subgroup of order 2 since this subgroup includes the identity and the generating element has
order 2 and is its own inverse. Adding another element of $Q_8$ and its inverse will result in one 
of the listed subgroups of order 4. We cannot add a single pair to a subgroup of order 4 since
the order of the resulting set will be 6, which is not a divisor of 8. Thus our list is exhaustive.

\noindent(c) Which of these subgroups is normal? For each one that is normal, name a group
isomorphic to that quotient.

All of the subgroups listed are normal. The quotient of the trivial subgroup $Q_8/\langle1\rangle$ is isomorphic to 
$Q_8$ and the quotient of the entire group $Q_8/Q_8$ is isomorphic to $\mathbb{Z}_1$.
The first listed subgroup is normal because multiplication with the identity and -1 in $Q_8$ is commutative.
Furthermore without loss of generality, we let subgroup $N = \{1, -1, i, -i\}$. $1 \cdot N \cdot -1 = N$ and for
$jNj^{-1}$ we check $jij^{-1} \in N$. Since $ji=k$ and $j^{-1} = -j$ we have $jij^{-1} = ji(-j) = (ji)(-j) = k(-j) = -(kj) = -(-i) = i$ since $ji=k$ and $j^{-1} = -j$.
Since $jij^{-1} = i$ we have that N is a normal subgroup of $Q_8$. This argument applies symmetrically for elements $j$ and $k$.
Thus we have shown that the listed elements of order 4 are normal subgroups.

\begin{itemize}
    \item $\{1, -1\} \cong \mathbb{Z}_2$
    \item $\{1, -1, i, -i\} \cong \mathbb{Z}_4$
    \item $\{1, -1, j, -j\} \cong \mathbb{Z}_4$
    \item $\{1, -1, k, -k\} \cong \mathbb{Z}_4$
\end{itemize}


\section*{Problem 4}
\noindent\textit{Let $M$ and $N$ be normal subgroups of a group $G$. Prove that
$MN = \{mn | m \in M \text{, } n \in N \}$ is a normal subgroup of $G$. }

\begin{proof}
    Given that $M$ and $N$ are both normal subgroups of $G$, then we know that
    $g^{-1}mg \in M$ and $g^{-1}ng \in N$.
    
    \begin{centering}
        $(g^{-1}mg)(g^{-1}ng) = g^{-1}m(gg^-1)ng = g^{-1}(mn)g \in MN \implies$ $MN$ is normal.
    \end{centering}
    Additionally we show that MN is a subgroup by showing closure and inverses.
    \begin{itemize}
        \item \textbf{Inverse} $mn \implies g^{-1}mng$ and $(g^{-1}mng)^{-1} = (g^{-1}n^{-1}m^{-1}g) \implies
        (mn)^{-1} \in MN$
        \item \textbf{Closure} Given $mn, m'n' \in MN \implies mn(m'n') = m(nm')n$. From the normality of
        $M$ we have that $n'm' = m^*n'$. Thus we have $m(nm')n = m(m^*n')n = (mm^*)(n'n) \in MN$
    \end{itemize}
    Thus since $MN$ is closed under the group operation and every element has an inverse, $MN$ is a subgroup.
\end{proof}

\section*{Problem 5}
\noindent\textit{Let $G$ be a group, and let Aut($G$) be the group of automorphisms of $G$. Recall that
an inner automorphism of $G$ is an automorphism of the form $\varphi_g(x) = gxg^{-1}$ for some
$g \in G$. Let Inn($G$) be a subset of Aut($G$) denote the set of inner automorphisms of $G$.}

\noindent\text{(a) Show that Inn($G$) $\cong G/Z(G)$, where $Z(G)$ is the center of $G$.}
\begin{proof}
    Let $\Phi : G \to \text{Inn}(G)$ where $g \mapsto \varphi_g$. For elements $g, h \in G$ we have
    $\Phi(gh)(x) = (gh)x(gh)^{-1} = g(hxh^{-1})g^{-1} = \Phi_g(\Phi_h(x)) = \Phi_g \circ \Phi_h(x)$
    and thus $\Phi$ is a group homomorphism.
    We define the kernal of $\Phi$. The element $g \in \text{Ker }\Phi$ if and only if 
    \begin{align*}
        \varphi_g(x) &= x &\quad \text{for all x} \\
        gxg^{-1} &= x &\quad \text{for all x} \\
        gx &= xg &\quad \text{for all x} \\
    \end{align*}
    thus $g \in Z(G)$ which is the kernal of $\Phi$.

    By the First Isomorphism Theorem, if we define Ker $\Phi$ as $Z(G)$ and Im $\Phi$ as
    Inn($G$), then we have the following:
    \begin{centering}
        $G/Z(G) \cong \text{Inn}(G)$
    \end{centering}
\end{proof}

\noindent\text{(b) Show that Inn$(G)$ is a normal subgroup of Aut$(G)$.}

\begin{proof}
    To show that Inn$(G)$ is a normal subgroup of Aut($G$), we check for all 
    $\psi \in \text{ Aut}(G)$ and $\varphi_g \in \text{Inn}(G)$ that $\psi \circ \varphi_g \circ \psi^{-1} \in \text{ Inn}(G)$
    which means it is an inner automorphism.
    \begin{align*}
        (\psi \circ \varphi_g \circ \psi^{-1})(x) &= \psi(\varphi_g(\psi^{-1}(x))) \\
        &= \psi(g\psi^{-1}(x)g^{-1}) \\
        &= \psi(g)\psi(\psi^{-1}(x))\psi(g^{-1}) \\
        &= \psi(g)x\psi(g)^{-1} \\
        &= \varphi_{\psi(g)}(x) \\
    \end{align*}
    We have shown that $(\psi \circ \varphi_g \circ \psi^{-1})(x) = \varphi_{\psi(g)}(x)$
    where $\varphi_{\psi(g)}$ is an inner automorphism defined by $\psi(g)$ and thus
    Inn$(G)$ is a normal subgroup of Aut($G$).
    
\end{proof}

\end{document}